% =====================================================================
%  FICHE 8 — THÈME 5 — NIVEAU DIFFICILE — ÉNONCÉS
% =====================================================================
\documentclass[11pt,a4paper]{article}
\usepackage[utf8]{inputenc}
\usepackage[T1]{fontenc}
\usepackage{lmodern}
\usepackage[french]{babel}
\usepackage{amsmath,amssymb,mathtools}
\usepackage{icomma}
\usepackage{siunitx}
\sisetup{output-decimal-marker={,},per-mode=power,group-separator={\,},group-minimum-digits=5}
\usepackage[version=4]{mhchem}
\usepackage{xcolor}
\usepackage{colortbl}
\usepackage{tikz}
\usepackage[most]{tcolorbox}
\usepackage{enumitem}
\usepackage{array}
\usepackage{booktabs}
\usepackage{fancyhdr}
\usepackage[hidelinks]{hyperref}
\usetikzlibrary{arrows.meta,positioning,calc}
\usepackage[a4paper,margin=2.1cm,headheight=26pt,headsep=10pt]{geometry}
\usepackage{titlesec}

\definecolor{primary}{RGB}{150,98,162}
\definecolor{primarydark}{RGB}{92,52,110}

\tcbset{boxbase/.style={enhanced,breakable,boxrule=0.9pt,arc=2.5pt,left=3mm,right=3mm,top=2.3mm,bottom=2.3mm,fonttitle=\bfseries,coltitle=white,toptitle=0.6mm,bottomtitle=0.6mm}}
\newtcolorbox[auto counter]{exercice}[1][]{boxbase,colback=primary!4,colframe=primary,title={\textsc{Exercice}~\thetcbcounter\ \ifx\relax#1\relax\relax\else\textmd{--- \itshape #1}\fi}}

\newcommand{\dHr}{\Delta H^{\circ}_{\text{r}}}

\pagestyle{fancy}
\fancyhf{}
\fancyhead[L]{\small\textcolor{primarydark}{\textbf{Fiche 8} — Thème 5 : Équilibre chimique et Le Chatelier}}
\fancyhead[R]{\small\textcolor{primarydark}{\textsc{Difficile}}}
\fancyfoot[C]{\small\thepage}
\renewcommand{\headrulewidth}{0.8pt}
\renewcommand{\headrule}{\hbox to\headwidth{\color{primary}\leaders\hrule height \headrulewidth\hfill}}
\setlength{\parindent}{0pt}\setlength{\parskip}{3pt}

\begin{document}

\begin{center}
\begin{tikzpicture}
\node[rectangle,rounded corners=7pt,fill=primary,text=white,minimum width=16cm,
      minimum height=1.1cm,align=center,inner sep=6pt]
  {{\large\bfseries Thème 5 — Équilibre chimique et Le Chatelier}\\[1pt]{\small Niveau \textsc{Difficile} — énoncés}};
\end{tikzpicture}
\end{center}

\begin{exercice}[tableau d'avancement et équation du second degré]
On introduit \SI{1.0}{\mol} de tétraoxyde de diazote dans une enceinte fermée de \SI{1.0}{\litre}, où
s'établit l'équilibre \;$\ce{N2O4}(g) \rightleftharpoons 2\,\ce{NO2}(g)$,\; de constante $K = 2{,}0$ (en
concentrations).
\begin{enumerate}[label=\alph*),leftmargin=2em,topsep=2pt,itemsep=2pt]
  \item En notant $x$ la concentration de $\ce{N2O4}$ ayant réagi, dresser le tableau des concentrations
        à l'équilibre.
  \item Écrire $K$ en fonction de $x$.
  \item Résoudre l'équation du second degré obtenue et en déduire $x$.
  \item Donner les concentrations de $\ce{N2O4}$ et de $\ce{NO2}$ à l'équilibre.
  \item Vérifier le résultat en recalculant $K$.
\end{enumerate}
\end{exercice}

\begin{exercice}[équilibre symétrique : l'astuce de la racine carrée]
On introduit \SI{2.0}{\mol} de $\ce{H2}$ et \SI{2.0}{\mol} de $\ce{I2}$ dans un réacteur de \SI{1.0}{\litre},
où s'établit \;$\ce{H2}(g) + \ce{I2}(g) \rightleftharpoons 2\,\ce{HI}(g)$,\; de constante $K = 16$.
\begin{enumerate}[label=\alph*),leftmargin=2em,topsep=2pt,itemsep=2pt]
  \item En notant $x$ la concentration ayant réagi, dresser le tableau des concentrations à l'équilibre.
  \item Écrire $K$ en fonction de $x$. Pourquoi peut-on ici prendre directement la \emph{racine carrée}
        des deux membres ?
  \item En déduire $x$.
  \item Donner la concentration de $\ce{HI}$ à l'équilibre.
\end{enumerate}
\end{exercice}

\begin{exercice}[du taux de décomposition à la constante $K$]
L'hydrogénosulfure d'ammonium se décompose selon
\;$\ce{NH4HS}(s) \rightleftharpoons \ce{NH3}(g) + \ce{H2S}(g)$.\; On introduit \SI{12.75}{\gram} de
$\ce{NH4HS}$ dans une enceinte de \SI{2.00}{\litre} à \SI{25}{\celsius}. À l'équilibre, \SI{20}{\percent}
du sel s'est décomposé. On donne $M(\ce{NH4HS}) = \SI{51}{\gram\per\mol}$.
\begin{enumerate}[label=\alph*),leftmargin=2em,topsep=2pt,itemsep=2pt]
  \item Calculer la quantité de matière initiale $n_0$ de $\ce{NH4HS}$.
  \item Quelle quantité s'est décomposée ? En déduire les quantités de $\ce{NH3}$ et de $\ce{H2S}$ formées.
  \item Calculer les concentrations de $\ce{NH3}$ et de $\ce{H2S}$ à l'équilibre.
  \item Écrire $K_c$ (attention au solide) et calculer sa valeur numérique.
\end{enumerate}
\end{exercice}

\begin{exercice}[optimiser un équilibre : tous les leviers réunis]
Dans un réacteur, on considère l'équilibre \;$\ce{C}(s) + \ce{CO2}(g) \rightleftharpoons 2\,\ce{CO}(g)$,\;
avec $\dHr = \SI{+172}{\kilo\joule\per\mol}$ (réaction \textbf{endothermique}). On cherche à
\textbf{augmenter le rendement en $\ce{CO}$}.
\begin{enumerate}[label=\alph*),leftmargin=2em,topsep=2pt,itemsep=2pt]
  \item Écrire la constante $K$ (ne pas oublier d'exclure le carbone solide).
  \item Compter les moles de gaz de chaque côté. Quel est l'effet d'une \emph{augmentation du volume}
        (donc d'une baisse de pression) sur l'équilibre ?
  \item Quel est l'effet d'une \emph{augmentation de la température} sur l'équilibre et sur la valeur de $K$ ?
  \item Quel est l'effet (i) d'un retrait de $\ce{CO2}$, (ii) de l'ajout d'un catalyseur ?
  \item Conclure : citer les opérations qui augmentent effectivement le rendement en $\ce{CO}$.
\end{enumerate}
\end{exercice}

\end{document}
